\chapter{Realisierung und Anwendung von Softwarearchitekturen 1}
\label{ch:02_implementation-and-application-of-software-architectures_01}

\begin{universityTask}
    \begin{enumerate}[label=(\alph*)]
        \item Das \q{Strategic Alignment Model} liefert Informationen über die grundsätzliche Absicht des IT-Einsatzes. Benennen Sie seine vier Perspektiven und formulieren Sie jeweils ein Beispiel aus der Unternehmenspraxis.
            \universityPoints{8}
        \item Ergänzen Sie in der nachfolgenden Integrationsdimensionsmatrix die Ausprägungen der fünf Integrationsdimensionen.
            \begin{table}[H]
                \centering
                \begin{university-table}{
                    width=0.8\textwidth,
                    colspec = {Q[l,wd=5cm] X X X X X X X X X X X X},
                }
                    Integrationsdimension   & \SetCell[c=12]{l} Ausprägung & & & & & & & & & & & \\
                    Gegenstand              & \SetCell[c=3]{} & & & \SetCell[c=3]{} & & & \SetCell[c=3]{} & & & \SetCell[c=3]{} & & \\
                    Richtung                & \SetCell[c=6]{} & & & & & & \SetCell[c=6]{} & & & & & \\
                    Reichweite              & \SetCell[c=6]{} & & & & & & \SetCell[c=6]{} & & & & & \\
                    Realisierung            & \SetCell[c=4]{} & & & & \SetCell[c=4]{} & & & & \SetCell[c=4]{} & & & \\
                    Automatisierungsgrad    & \SetCell[c=6]{} & & & & & & \SetCell[c=6]{} & & & & & \\
                \end{university-table}
            \end{table}
            \universityPoints{10}
        \item Bearbeiten Sie die nachfolgende Fallstudie zur Organisationsmodellierung. Definieren Sie zunächst die drei beteiligten Akteure und erstellen Sie ein grobes Interaktionsmodell, das die Interaktionen modellierungsmethodisch abbildet. Verfeinern Sie daraufhin in einem zweiten Schritt das zunächst erstellte Interaktionsmodell, indem Sie die einzelnen Organisationseinheiten der Unternehmung in Ihre Modellierung einbeziehen. Als Organisationseinheiten verwenden Sie den \q{Vertrieb}, die \q{Fertigung}, die \q{Beschaffung} und den \q{Versand}. \linebreak
            \begin{quotation}
                \begin{itshape}
                    Der Kunde Dr. Müller bestellt bei der Unternehmung Gartenschlauch.de einen speziellen Gartenschlauch, der aufgrund seiner Sondergröße zunächst angefertigt werden muss. Die Bestellung wird zunächst auf ihre Machbarkeit hin geprüft. Nach Auftragsannahme werden die benötigten Materialien bei einem Lieferanten beschafft. Nach Eintreffen des Materials und Einplanung des Auftrags werden die Artikel anhand eines Arbeitsplans gefertigt und an Dr. Müller verschickt.
                \end{itshape}
            \end{quotation}
            \universityPoints{12}
    \end{enumerate}
\end{universityTask}

\section{Strategic Alignment Model}
\label{sec:02-01_strategic-alignment-model}

Der Text enthält die Lösung der obigen Aufgabenstellung.

\section{Integrationsdimensionen}
\label{sec:02-02_integration-dimensions}

Der Text enthält die Lösung der obigen Aufgabenstellung.

\section{Organisationsmodellierung}
\label{sec:02-03_organizational-modeling}

Der Text enthält die Lösung der obigen Aufgabenstellung.